\thispagestyle{normalpage}

\begin{center}
    \large \textbf{CHAPITRE I} \\[0.3em]
    \large \textbf{CONTEXTE GÉNÉRAL DU PROJET}
\end{center}
\addcontentsline{toc}{chapter}{CHAPITRE I - CONTEXTE GÉNÉRAL DU PROJET}
\label{page:chapitre1}
\section{I.1. Présentation de l'INSI}
\label{sec:insi}

\subsection{I.1.1. Information d'ordre général}

INSI (Institut National Supérieur d'Informatique) est un institut universitaire situé à Antananarivo, spécialisé dans le domaine de l'informatique. Il accueille principalement des bacheliers ayant choisi cette filière, mais reste également ouvert à toute personne manifestant un intérêt particulier pour les métiers du numérique, à travers son centre de formation professionnelle. L'établissement a été fondé le 17 août 2021 et son siège se trouve à Ankazotokana Ambanidia, Lot VF 32 Ter.

Pour tout renseignement, vous pouvez nous contacter par email à contact@insi.mg ou visiter notre site web : www.insi.mg

INSI University bénéficie du soutien d'entreprises partenaires, dont \textbf{Ilo Madagascar SARLU} et \textbf{Data Pulse Center}, qui sont également à l'origine de sa création.

\subsection{I.1.2. Missions et historiques}

INSI University a été fondée le 17 août 2021 à Antananarivo, dans un contexte marqué par une forte demande de compétences dans le domaine de l'informatique et des technologies numériques.

Créé à l'initiative de \textbf{Ilo Madagascar SARLU} et \textbf{Data Pulse Center}, l'Institut a vu le jour avec pour ambition de former une nouvelle génération de professionnels qualifiés, capables de répondre aux besoins du marché local et international.

Depuis sa création, INSI ne cesse d'évoluer en adaptant ses formations aux enjeux technologiques actuels, tout en favorisant un apprentissage pratique, orienté vers les réalités du monde professionnel.

La mission principale de l'INSI est de \textbf{former des experts compétents, éthiques et innovants dans le domaine de l'informatique}, capables de contribuer activement au développement économique et numérique de Madagascar.

L'Institut s'engage à :

\begin{itemize}
    \item Offrir des formations de qualité, à jour des évolutions technologiques,
    \item Favoriser l'insertion professionnelle de ses apprenants,
    \item Créer un lien fort entre la formation académique et les besoins réels des entreprises,
    \item Encourager l'entrepreneuriat et l'innovation technologique,
    \item Contribuer à la démocratisation de l'accès au savoir numérique.
\end{itemize}

\subsection{I.1.3. Organigramme de l'INSI}

L'organigramme ci-dessous présente la structure organisationnelle de l'INSI, mettant en évidence les différentes fonctions et responsabilités au sein de l'institut, ainsi que les liens hiérarchiques entre les divers services.

\label{page:f1}
\begin{figure}[H]
    \centering
    % \includegraphics[width=0.9\textwidth]{usecase.png} 
    \caption{Organigramme de l'INSI}
    \vspace{0.2em}
    \small (Source : INSI, 2025)
    \label{fig:organigramme_insi}
\end{figure}

L'organigramme de l'INSI, en partenariat avec le Collège de Paris, est structuré autour d'une direction générale incarnée par le Président Directeur Général. Ce dernier assure la coordination de
trois pôles principaux : la \textbf{Direction Administrative et Financière}, la \textbf{Direction des Études}, et le service des \textbf{Relations Internationales et Mobilités}.

Le \textbf{Président Directeur Général} assure la direction stratégique de l'institut. Il supervise l'ensemble des activités pédagogiques, administratives et de développement.

Il veille à la qualité de l'enseignement, à la bonne gestion des ressources humaines et matérielles, et représente l'établissement auprès des partenaires et autorités.

La \textbf{Direction Administrative et Financière} est chargée de la gestion des affaires internes, notamment à travers deux services essentiels : le \textbf{service des scolarités}, qui s'occupe des dossiers étudiants, inscriptions et suivi académique, et les \textbf{services supports et stratégiques}, qui assurent un appui logistique, administratif et technologique à l'ensemble de l'institut.

Le pôle des \textbf{Relations Internationales et Mobilités} gère les échanges académiques avec l'étranger ainsi que les projets de mobilité étudiante ou enseignante.

La \textbf{Direction des Études} supervise quant à elle l'ensemble des programmes pédagogiques et le corps enseignant. Elle travaille en étroite collaboration avec un \textbf{Responsable Parcours}, qui coordonne les différentes spécialisations proposées. Cette direction intègre également une cellule dédiée à la \textbf{formation professionnelle et à la certification}, en charge des dispositifs de validation des compétences, des certifications professionnelles et des parcours de formation continue. Il comprend également le \textbf{Département alternance et emplois}, responsable du suivi des étudiants en contrat d'alternance, des stages professionnels, et de l'insertion professionnelle.

Sous la partie \textbf{Responsable Parcours}, plusieurs responsables pédagogiques gèrent les différentes filières de formation : le \textbf{parcours tronc commun}, qui regroupe les matières fondamentales partagées par tous les étudiants en début de ciclo ; le \textbf{parcours ARSB \& IC} ; le \textbf{parcours GL \& MSI} ; et enfin le \textbf{parcours IAD \& IAD}.

D'autres structures, bien que non représentés sur cet organigramme, contribuent également de manière significative à la mission de l'INSI.

Le \textbf{Conseil de l'École}, il regroupe des membres de l'équipe pédagogique et administrative. Il joue un rôle consultatif ou décisionnel sur les questions académiques : validation des programmes, suivi des résultats, discipline, méthodologie, etc.

Le \textbf{responsable de la formation en ligne}, il conçoit et supervise les dispositifs d'apprentissage à distance. Il s'assure que les contenus en ligne sont accessibles, interactifs et conformes aux standards pédagogiques, tout en accompagnant enseignants et étudiants dans leur utilisation.

L'\textbf{Association et Clubs étudiants}, encadré et soutient les initiatives étudiantes, les clubs thématiques (informatique, innovation, sport, culture, etc.), et les événements organisés dans le

cadre de la vie étudiante. Cela favorise l'engagement, le développement personnel et les compétences transversales des étudiants.

Cette structure témoigne d'une organisation à la fois pédagogique, administrative et tournée vers l'international et l'employabilité, assurant une formation complète et professionnalisante aux étudiants.

\subsection{I.1.4. Domaine de spécialisation}

INSI University propose des \textbf{parcours académiques professionnalisants} dans le domaine de l'informatique et des technologies numériques, articulés autour de licences professionnelles, de masters spécialisés, formations professionnelles continues et d'opportunités de mobilité internationale.

\subsubsection*{a) Licences professionnelles}

Les programmes de licence offrent une formation de base solide, axée sur la pratique, et préparant à une insertion rapide dans le monde professionnel :



\begin{spacing}{1.05} % interligne un peu réduit

\begin{itemize}[noitemsep,nolistsep]
    \item \textbf{GL (Génie Logiciel)}: Conception, développement et maintenance d'applications logicielles.
    \item \textbf{ARSB (Administration Réseaux et Systèmes et Base de données)} : Gestion d'infrastructures informatiques, réseaux, serveurs et systèmes.
    \item \textbf{IAD (Intelligence Artificielle et Data Science)} : Introduction aux algorithmes d'IA, au traitement de données et à l'apprentissage automatique.
\end{itemize}

\subsubsection*{b) Masters professionnels}

Les masters permettent d'approfondir les compétences techniques et managériales, en lien étroit avec les évolutions du secteur :

\begin{itemize}[noitemsep,nolistsep]
    \item \textbf{MSI (Management des Systèmes d'Information)} : Pilotage de projets numériques, gouvernance des SI.
    \item \textbf{IC (Infrastructures et Cybersécurités)} : Sécurisation des systèmes, audit, réseaux avancés.
    \item \textbf{I2AD (Intelligence Artificielle, Analyse de Données et Automatisation)} : Applications avancées de l'IA, automatisation des processus et science des données.
\end{itemize}

\subsubsection*{c) Formations professionnelles}

INSI University propose également des formations professionnelles continues, spécifiquement conçues pour les professionnels déjà en poste souhaitant approfondir leurs compétences ou se spécialiser dans l'un des domaines proposés. Ces formations sont adaptées aux contraintes des professionnels (horaires flexibles, modules pratiques) et visent à renforcer leur employabilité ou à accompagner leur évolution de carrière.

\subsubsection*{d) Mobilité internationale \& double diplôme}

Dans une optique d'ouverture et de reconnaissance internationale, INSI propose à ses étudiants des opportunités de mobilité académique et de double diplôme, à partir de la deuxième année (L2).

\begin{itemize}[noitemsep,nolistsep]
    \item \textbf{Format hybride} : 6 mois de formation à Madagascar + 6 mois à l'étranger selon le parcours choisi.
    \item \textbf{Double diplôme en L3 ou M2} : En partenariat avec des institutions internationales telles que :
    \begin{itemize}[noitemsep,nolistsep]
        \item Collège de Paris (France)
        \item Keyce Informatique (France)
        \item Eskills (Québec, Canada)
    \end{itemize}
\end{itemize}

%\vspace{0.5cm} % supprimé pour gagner de la place

\subsubsection*{Conditions d'éligibilité}

\begin{itemize}[noitemsep,nolistsep]
    \item Niveau DALF C1 en français pour la mobilité vers des pays francophones.
    \item Niveau B2 ou C1 en anglais pour les destinations anglophones.
\end{itemize}

\subsection{I.1.5 Architectures des formations pédagogiques}

\noindent
{\footnotesize
\begin{spacing}{0.9}
\begin{minipage}{0.95\textwidth}
Au sein de l'INSI, une seule mention est proposée : Informatique, déclinée en trois parcours distincts aux niveaux Licence et Master :
\end{minipage}
\end{spacing}
}

\end{spacing}

 % fin \small

% fin \small
% Dans le préambule, ajoute (si pas déjà présent) :

% Dans le document :

\newpage
\label{page:tableau1}
\begin{table}[H]
    \centering
    \begin{tabular}{|p{6cm}|p{6cm}|}
        \hline
        \textbf{Licence} & \textbf{Master} \\
        \hline
        Génie Logiciel (GL) & Management des Systèmes d'Information (MSI) \\
        \hline
        Administration Réseaux, Systèmes et Bases de données (ARSB) & Infrastructures et Cybersécurité (IC) \\
        \hline
        Intelligence Artificielle et Data Science (IAD) & Intelligence Artificielle, Analyse de Données et Automatisation (IADA) \\
        \hline
    \end{tabular}
    \caption{Les parcours en Licence et Master}
    \label{tab:parcours_licence_master}
    \small\vspace{0.3em}
    \small (Source : INSI,2025)
\end{table}

{\small % texte en taille réduite

\setstretch{1.0} % interligne simple

L'architecture des études à trois niveaux conformément au système LMD (Licence Master et Doctorat) permet les comparaisons et les équivalences académiques des diplômes au niveau international.

L = Licence (Bac + 3) = L1, L2, L3 = 6 semestres S1 à S6

M = Master (Bac + 5) = M1, M2 = 4 semestres S7 à S10

Le diplôme de licence est obtenu en 3 années des études après Baccalauréat. Et le diplôme de Master est obtenu en 2 ans après obtenu du diplôme de LICENCE.

Le MASTER PROFESSIONNEL est un diplôme destiné à la recherche emploi au terme des études.

\subsubsection*{Spécificités pédagogiques}

\begin{itemize}[noitemsep,nolistsep,left=0pt]
    \item École d'alternance dès la L2 S4 : alternance cours/entreprise pour une meilleure insertion
    \item Écoles connectées \& e-learning : plateforme numérique de formation, cloud intégré, Google Workspace. Cours accessibles en ligne selon les parcours
    \item Pédagogie active \& par projets : hackathons, soutenances exposées, projets réels en entreprise
\end{itemize}

\subsection{I.1.6. Relations de l'INSI avec les entreprises et les organismes}

L'Institut National Supérieur en Informatique (INSI) entretient des liens étroits avec le monde professionnel. L'INSI collabore activement avec de nombreux organismes publics, semi-publics et privés, tant au niveau national qu'international. Cette collaboration se traduit concrètement par l'organisation régulière de stages obligatoires pour les étudiants, favorisant leur immersion professionnelle et leur employabilité.

Parmi les entreprises et institutions partenaires accueillant nos étudiants en stage chaque année, on peut citer : Data Pulse Center, RELIA Consulting, Code Talent, Bocasay, Minimad, Smartone AI, NOVITY, Eclosia Madagascar, VISEO Groupe, Océane Trade,\\
ACCES Banque,

} % fin small


\subsection{I.1.7. Partenariat au niveau international}

L'INSI s'est inscrit dans une dynamique d'ouverture à l'international grâce à son partenaire fondateur, \textbf{ILO Madagascar}, représentant de l'Organisation Internationale du Travail. Cette orientation se renforce aujourd'hui à travers des \textbf{partenariats stratégiques avec des institutions académiques et des entreprises internationales}, permettant à l'Institut de proposer une formation alignée sur les standards mondiaux.

L'INSI bénéficie actuellement de \textbf{programmes de co-diplomation et de double diplôme} en collaboration avec plusieurs établissements internationaux de renom, tels que :

\begin{itemize}
    \item \textbf{Collège de Paris International}
    \item \textbf{Keyce Informatique} (France)
    \item \textbf{Eskills Academy} (Québec, Canada)
    \item \textbf{École Supérieure Polytechnique d'Amistranana (ESPA)} -- pour une synergie régionale renforcée.
\end{itemize}

Ces partenariats permettent non seulement d'enrichir l'offre pédagogique, mais aussi d'ouvrir des perspectives de \textbf{mobilité académique et professionnelle} aux étudiants, en leur offrant une reconnaissance de leurs compétences à l'échelle internationale.

Par ailleurs, à travers les \textbf{stages annuels et projets pédagogiques}, l'INSI collabore également avec des entreprises à dimension régionale ou mondiale, telles que :

\begin{itemize}
    \item \textbf{VISEO Groupe}
    \item \textbf{Edoisia Madagascar} (Groupe mauricien)
    \item \textbf{Smartone A1}
    \item \textbf{NOVITY}, entre autres.
\end{itemize}

Ces relations actives témoignent de l'ancrage de l'INSI dans un \textbf{écosystème global}, plaçant l'Institut comme un acteur majeur de la formation en informatique à Madagascar, tourné vers l'innovation et les exigences du marché international.

\subsection{I.1.8. Parcours professionnels des diplômés}

Les diplômés de la Licence en Informatique de l'INSI disposent d'un socle solide de compétences en programmation, systèmes, réseaux, bases de données et développement logiciel. Cette polyvalence leur permet d'intégrer directement le marché du travail ou de poursuivre en Master. À ce niveau, ils peuvent exercer en tant que :

\begin{itemize}
    \item \textbf{Technicien système et réseau}
    \item \textbf{Développeur web / mobile (front-end ou back-end)}
    \item \textbf{Administrateur système junior}
    \item \textbf{Technicien support IT}
    \item \textbf{Intégrateur / testeur logiciel}
    \item \textbf{Assistant en sécurité informatique}
    \item \textbf{Assistant en projets d'IA ou de data (selon profil)}
\end{itemize}

Les stages annuels en entreprise leur offrent une première expérience professionnelle valorisante dans des structures variées : entreprises technologiques, banques, administrations, startups, ONG, etc.

La formation de Master proposée à l'INSI, grâce à ses parcours spécialisés, forme des professionnels capables d'intervenir sur des projets complexes, innovants et à dimension stratégique. Les débouchés varient selon la spécialisation choisie :

\subsubsection*{Parcours Ingénierie logicielle, réseaux et systèmes :}
\begin{itemize}
    \item \textbf{Ingénieur en développement logiciel / full-stack}
    \item \textbf{Architecte logiciel}
    \item \textbf{Ingénieur système et réseaux}
    \item \textbf{Administrateur cloud / DevOps}
    \item \textbf{Chef de projet informatique}
\end{itemize}

\subsubsection*{Parcours Cybersécurité :}
\begin{itemize}
    \item \textbf{Expert en sécurité des systèmes d'information}
    \item \textbf{Analyste SOC / CERT}
    \item \textbf{Auditeur en cybersécurité}
    \item \textbf{Responsable de la sécurité IT (RSSI)}
\end{itemize}

\subsubsection*{Parcours Intelligence Artificielle et Data :}
\begin{itemize}
    \item \textbf{Data Scientist}
    \item \textbf{Data Analyst}
    \item \textbf{Ingénieur en Intelligence Artificielle}
    \item \textbf{Machine Learning Engineer}
\end{itemize}

\begin{itemize}
    \item \textbf{Développeur d'applications intelligentes}
    \item \textbf{Consultant en IA}
    \item \textbf{Spécialiste en traitement automatique du langage naturel (TALN)}
    \item \textbf{Spécialiste en vision par ordinateur (Computer Vision)}
\end{itemize}

Grâce aux partenariats internationaux et aux projets en entreprise, les diplômés du Master peuvent également prétendre à des postes dans des structures multinationales, ou poursuivre une carrière académique et de recherche (doctorat, enseignement supérieur).

\subsection{I.1.9. Ressources humaines}

L'INSI s'appuie sur une équipe pédagogique et administrative compétente, dynamique et engagée dans la réussite de ses étudiants. Ses ressources humaines se répartissent comme suit :

\begin{itemize}
    \item Président Directeur de l'INSI : Dr RAMASONDRANO Andriamanjaka
    \item Responsable parcours ARSB \& IC : Mr RATIANARISON Hery Mampionona
    \item Responsable parcours tronc commun : Dr RAKOTONDRAMANANA R, Sitraka
    \item Responsable parcours GL \& MSI : Mr ANDRIAMIADANA TANTELY ANTOINE
    \item Responsable parcours IAD \& I2AD : par intérim PDG
    \item Nombre des enseignants : 40
    \item Nombre des personnels administratives : 04
    \item Personnels d'appui : 6
\end{itemize}

\section{I.2. Contexte professionnel et entreprise d'accueil}
\subsection{I.2.1. Présentation de l'entreprise d'accueil}
L'entreprise d'accueil de ce projet est \textbf{Techcare}, une startup technologique specialisée dans la vente 
de produits pharmaceutiques en ligne à Madagascar.\\ Fondée en 2024, Techcare vise à révolutionner l'accès aux médicaments et aux produits de santé en offrant une plateforme e-commerce conviviale et fiable.\\
Techcare collabore avec un réseau de pharmacies partenaires à travers le pays pour garantir la disponibilité
des produits et assurer une livraison rapide aux clients.\\
L'entreprise met l'accent sur la qualité du service client, la sécurité des transactions en ligne, et la conformité aux réglementations locales en matière de santé et de commerce électronique.\\
Grâce à son approche innovante, Techcare aspire à devenir le leader du marché de la pharmacie en ligne à Madagascar, en facilitant l'accès aux soins de santé pour tous les citoyens, même dans les régions éloignées.


\subsection{I.2.2. Contexte du projet}
Le projet de migration d'une architecture monolithique vers une architecture basée sur des microservices s'inscrit dans le cadre de l'évolution technologique de Techcare.\\
Avec la croissance rapide du commerce en ligne et l'augmentation du nombre d'utilisateurs, 
l'architecture monolithique actuelle de la plateforme e-commerce de Techcare présente des limitations en termes de scalabilité, de maintenabilité et de flexibilité.\\
La décision de migrer vers une architecture basée sur des microservices vise à répondre à ces défis en permettant une meilleure gestion des ressources, une adaptation plus rapide aux besoins changeants du marché, et une amélioration de l'expérience utilisateur.\\
Ce projet est également motivé par la volonté de Techcare d'adopter des pratiques de développement logiciel modernes, favorisant l'innovation et la compétitivité dans un secteur en constante évolution.

\section{I.3 Problématique et justification du projet}
La problématique centrale de ce projet est de déterminer comment migrer efficacement une architecture monolithique vers une architecture basée sur des microservices, tout en garantissant une expérience utilisateur fluide et cohérente.
Cette migration soulève plusieurs défis techniques, organisationnels et stratégiques, notamment en ce qui concerne la décomposition de l'application monolithique en services indépendants, la gestion des communications entre ces services, 
et la mise en place de processus de développement adaptés.\\
La justification de ce projet réside sur ces avantages clés offerts par les microservices :
\begin{itemize}
    \item Amélioration de la scalabilité: chaque microservice peut être mis à l'échelle indépendamment en fonction de la demande.
    \item Meilleure maintenabilité: les équipes peuvent travailler sur des services spécifiques sans affecter l'ensemble de l'application.
    \item Flexibilité accrue: possibilité d'adopter différentes technologies et langages pour différents services.
    \item Résilience améliorée: une défaillance dans un microservice n'affecte pas nécessairement l'ensemble du système.
\end{itemize}

\section{I.4 Objectifs et périmètre du projet}
L'objectif principal de ce projet est de migrer l'architecture monolithique existante de la plateforme e-commerce de Techcare vers une architecture basée sur des microservices, en utilisant des méthodologies et des techniques modernes de développement logiciel.\\
Les objectifs spécifiques incluent :
\begin{itemize}
    \item Analyser l'architecture monolithique actuelle et identifier les composants clés à migrer.
    \item Concevoir une architecture cible basée sur des microservices, en définissant les services, leurs interactions et les technologies à utiliser.
    \item Mettre en œuvre la migration en développant et déployant les microservices.
    \item Tester et valider la nouvelle architecture pour assurer sa performance, sa scalabilité et sa fiabilité.
    \item Documenter le processus de migration et les leçons apprises pour faciliter les futures évolutions.
\end{itemize}
la gestion des utilisateurs, l'inventaire et les notifications.

\section{I.5 Méthodologie et plan de travail}

Pour mener à bien ce projet de migration complexe, nous avons adopté une démarche structurée combinant des principes agiles et une planification rigoureuse par phases.

\subsection{I.5.1 Méthodologie de gestion de projet : Scrum}

Le choix s'est porté sur la méthodologie \textbf{Scrum}, une approche agile permettant de gérer le développement de manière itérative et incrémentale. Cette méthode est particulièrement adaptée à une architecture microservices où chaque service peut être considéré comme une unité fonctionnelle indépendante.

\begin{itemize}
    \item \textbf{Sprints} : Le travail a été découpé en cycles de deux semaines (Sprints). Chaque sprint se termine par la livraison d'une fonctionnalité ou d'un service testable.
    \item \textbf{Cérémonies} : Des réunions de planification (Sprint Planning) et des revues de sprint ont permis d'ajuster les priorités en fonction des défis techniques rencontrés (notamment sur la gestion des communications asynchrones).
    \item \textbf{Backlog} : Les besoins ont été traduits en "User Stories" regroupées dans un backlog produit, priorisées selon leur importance critique pour le système (commençant par l'authentification et la passerelle API).
\end{itemize}

\subsection{I.5.2 Plan de travail et étapes de réalisation}

Le projet s'articule autour de cinq phases majeures réparties sur la durée du stage :

\begin{enumerate}
    \item \textbf{Phase d'Analyse (Semaine 1-2)} : Étude approfondie du monolithe existant, identification des domaines fonctionnels (Bounded Contexts) et définition de la stratégie de découpage.
    \item \textbf{Phase de Conception (Semaine 3-4)} : Modélisation des microservices, choix des protocoles de communication (RabbitMQ pour l'asynchronisme, REST pour le synchrone) et conception des schémas de données distribués.
    \item \textbf{Phase d'Implémentation (Semaine 5-12)} : Développement itératif des microservices (User, Product, Inventory, Notification, File Manager) et de l'API Gateway.
    \item \textbf{Phase de Tests et Intégration (Semaine 13-14)} : Tests unitaires, tests d'intégration entre services et validation de la sécurité via JWT.
    \item \textbf{Phase de Déploiement et Documentation (Semaine 15-16)} : Conteneurisation des services, mise en place de l'infrastructure de production et finalisation du rapport technique.
\end{enumerate}

\subsection{I.5.3 Outils de collaboration et environnement technique}

La traçabilité et la qualité du code ont été assurées par l'utilisation d'outils standards de l'industrie :
\begin{itemize}
    \item \textbf{Git} : Pour la gestion de version et le travail collaboratif.
    \item \textbf{Concurrently} : Pour l'orchestration du démarrage des services en environnement de développement.
    \item \textbf{Docker} : Pour la gestion des dépendances d'infrastructure (PostgreSQL, MongoDB, RabbitMQ).
    \item \textbf{Trello/Jira} : Pour le suivi du tableau Kanban et l'avancement des tâches Scrum.
\end{itemize}

\section*{Conclusion partielle}
Ce premier chapitre a permis de poser les bases du projet en présentant l'organisme d'accueil, l'institut de formation et surtout les enjeux majeurs liés à la modernisation de la plateforme Techcare. L'analyse des besoins a mis en lumière les limites critiques de l'architecture monolithique actuelle face aux ambitions de croissance de l'entreprise. Cette compréhension du contexte justifie pleinement le passage vers une architecture distribuée, dont la conception détaillée fera l'objet du chapitre suivant.
